% generated by roam.py handbook export -- не редактировать, править заметки
% 20261025183020-метод_тяжелого_шарика.org (body 98992cece584)

\begin{method}[{Метод тяжёлого шарика}]\label[method]{note:C6A43F85-13E0-46B3-96A8-BCB9F34ECD57}
  Метод с моментом (инерцией):
  \[
  x_{k+1} = x_k - \alpha \nabla f(x_k) + \beta (x_k - x_{k-1})
  \]
  где $\alpha > 0$~--- шаг, $\beta \in [0;1)$~--- коэффициент момента; нужны две начальные точки $x_0, x_1$. Слагаемое $\beta(x_k - x_{k-1})$~--- инерция шарика, катящегося по графику: он продолжает двигаться в прежнем направлении.
\end{method}

\begin{remark}
  В отличие от \hyperref[note:BB5E3FC2-069C-44F0-92D0-B4DC13092CCA]{метода Нестерова} в форме без $v_k$, градиент считается в самой точке $x_k$, а не в точке после инерционного сдвига.
\end{remark}

\paragraph*{Теорема Поляка (без доказательства)}
Пусть $f \in C^2$, $x^{*}$~--- невырожденная точка локального минимума: $\mu E \le \nabla^2 f(x^{*}) \le LE$, $\mu > 0$. Пусть $0 \le \beta < 1$, $0 < \alpha < \frac{2(1+\beta)}{L}$.

Тогда найдутся $q \in [0;1)$ и $\varepsilon > 0$ такие, что для любых $x_0, x_1$ с $\|x_0 - x^{*}\| < \varepsilon$, $\|x_1 - x^{*}\| < \varepsilon$ метод сходится к $x^{*}$ и
\[
\forall \delta \in (0; 1-q)\ \exists C(\delta):\ \|x_k - x^{*}\| \le C(\delta)(q + \delta)^k
\]
Наименьшее $q$ достигается при
\[
\alpha^{*} = \frac{4}{(\sqrt{L} + \sqrt{\mu})^2},\quad \beta^{*} = \left( \frac{\sqrt{L} - \sqrt{\mu}}{\sqrt{L} + \sqrt{\mu}} \right)^2,\quad q^{*} = \frac{\sqrt{L} - \sqrt{\mu}}{\sqrt{L} + \sqrt{\mu}}
\]

\begin{remark}
  $q^{*} = \frac{1 - \sqrt{\mu/L}}{1 + \sqrt{\mu/L}} \approx 1 - 2\sqrt{\frac{\mu}{L}}$~--- та же зависимость от $\sqrt{\frac{\mu}{L}}$, что в \hyperref[note:C038105E-BD42-44A8-9542-2DBE9FFAAA3F]{сходимости метода Нестерова}, и лучше, чем $1 - \frac{\mu}{L}$ у \hyperref[note:85D9D77D-13A2-42E0-BF93-AD085E0134B9]{градиентного спуска}. Но гарантия только локальная.
\end{remark}
